\section{Семинар 2. Множества, логика и характеристические функции}
\noeqref{eq:quant_neg_all, eq:quant_neg_ex, eq:lim_def, eq:lim_neg}
\noeqref{eq:cap_cup_sub, eq:demorgan_cup, eq:char_func_def, eq:char_sub}
\noeqref{eq:bij_z, eq:bij_tan, eq:bij_hilbert}

\begin{remark}[Обзор темы]
    На втором семинаре разбираются правила построения отрицаний 
    кванторных высказываний, методы доказательства теоретико-множественных 
    тождеств (метод двух включений и алгебраический метод 
    характеристических функций), а также построение явных биекций для 
    доказательства равномощности бесконечных множеств.
\end{remark}

\subsection{Математическая логика и отрицание кванторов}

\begin{definition}
    В математической логике для формализации высказываний о множествах 
    используются два основных квантора:
    \begin{itemize}
        \item \textdef{Квантор всеобщности} $\forall$ («для любого», 
        «для каждого»);
        \item \textdef{Квантор существования} $\exists$ («существует хотя 
        бы один», «найдётся»).
    \end{itemize}
\end{definition}

\begin{properties}[Правила де Моргана для кванторов]
    При построении логического отрицания кванторного высказывания 
    каждый квантор заменяется на двойственный ему ($\forall \leftrightarrow 
    \exists$), а отрицание переносится на предикат:
    \begin{equation} \label{eq:quant_neg_all}
        \neg \bigl(\forall x \, P(x)\bigr) \iff \exists x \, \neg P(x) \tag{1}
    \end{equation}
    \begin{equation} \label{eq:quant_neg_ex}
        \neg \bigl(\exists x \, P(x)\bigr) \iff \forall x \, \neg P(x) \tag{2}
    \end{equation}
\end{properties}

\begin{example} ~
    Рассмотрим классическое определение предела числовой 
    последовательности $\{a_n\}$ на языке $\eps$-$N$. Число $A$ называется 
    пределом последовательности ($A = \lim_{n \to \infty} a_n$), если:
    \begin{equation} \label{eq:lim_def}
        \forall \eps > 0 \ \exists N \in \N \ \forall n > N : 
        \abs{a_n - A} < \eps \tag{3}
    \end{equation}
    Построим по правилам \eqref{eq:quant_neg_all} и \eqref{eq:quant_neg_ex} 
    логическое отрицание утверждения \eqref{eq:lim_def} (высказывание 
    «число $A$ \textbf{не} является пределом последовательности $\{a_n\}$»). 
    Заменяя все кванторы слева направо на противоположные и инвертируя 
    финальное неравенство, получаем:
    \begin{equation} \label{eq:lim_neg}
        \exists \eps > 0 \ \forall N \in \N \ \exists n > N : 
        \abs{a_n - A} \ge \eps \tag{4}
    \end{equation}
\end{example}

\subsection{Свойства включения и операции над множествами}

\begin{definition}
    Пусть $A$ и $B$ --- произвольные множества. Запишем основные 
    теоретико-множественные операции через логические связки:
    \begin{enumerate}
        \item \textdef{Объединение:} 
        $A \cup B = \{x \mid (x \in A) \lor (x \in B)\}$;
        \item \textdef{Пересечение:} 
        $A \cap B = \{x \mid (x \in A) \land (x \in B)\}$;
        \item \textdef{Разность:} 
        $A \setminus B = \{x \mid (x \in A) \land (x \notin B)\}$;
        \item \textdef{Симметрическая разность:} 
        $A \symdiff B = (A \setminus B) \cup (B \setminus A) = 
        \{x \mid (x \in A) \oplus (x \in B)\}$ (множество элементов, 
        принадлежащих ровно одному из двух множеств).
    \end{enumerate}
\end{definition}

\begin{problem}
    Доказать транзитивность отношения нестрогого включения множеств: 
    если $A \subseteq B$ и $B \subseteq C$, то $A \subseteq C$.
\end{problem}

\begin{proof}[Решение]~\\
    \begin{quote}
        По определению подмножества условие $A \subseteq B$ означает 
        $\forall x \, (x \in A \implies x \in B)$, а условие $B \subseteq C$ 
        означает $\forall x \, (x \in B \implies x \in C)$.

        Возьмём произвольный элемент $x \in A$. В силу включения 
        $A \subseteq B$ для него выполняется $x \in B$. Тогда в силу 
        включения $B \subseteq C$ из $x \in B$ следует $x \in C$. 
        Таким образом, $\forall x \, (x \in A \implies x \in C)$, что по 
        определению означает $A \subseteq C$.
    \end{quote}
\end{proof}

\begin{remark}
    Для любых двух множеств $A$ и $B$ справедлива цепочка включений:
    \begin{equation} \label{eq:cap_cup_sub}
        A \cap B \subseteq A \subseteq A \cup B \tag{5}
    \end{equation}
\end{remark}

\begin{proposition}[\textbf{Критерии включения множеств}]
    Для любых множеств $X$ и $Y$ следующие четыре утверждения попарно 
    эквивалентны:
    \[
        \text{1) } X \subseteq Y, \qquad
        \text{2) } X \setminus Y = \varnothing, \qquad
        \text{3) } X \cap Y = X, \qquad
        \text{4) } X \cup Y = Y.
    \]
\end{proposition}

\begin{proof}[Доказательство]~\\
    \begin{quote}
        Докажем эквивалентность утверждений по замкнутому циклу импликаций 
        $1 \implies 2 \implies 3 \implies 4 \implies 1$.

        \textbf{1) $\implies$ 2):} Пусть $X \subseteq Y$. Предположим от 
        противного, что $X \setminus Y \neq \varnothing$. Тогда существует 
        элемент $x \in X \setminus Y$, то есть $x \in X$ и $x \notin Y$. 
        Но из условия $X \subseteq Y$ и $x \in X$ должно следовать 
        $x \in Y$. Противоречие, значит, $X \setminus Y = \varnothing$.

        \textbf{2) $\implies$ 3):} Пусть $X \setminus Y = \varnothing$. 
        Включение $X \cap Y \subseteq X$ верно всегда согласно 
        \eqref{eq:cap_cup_sub}. Обратно, пусть $x \in X$. Если бы 
        выполнялось $x \notin Y$, то элемент $x$ принадлежал бы разности 
        $X \setminus Y$, которая пуста. Значит, $x \in Y$, откуда 
        $x \in X \cap Y$ и $X \subseteq X \cap Y$. Следовательно, 
        $X \cap Y = X$.

        \textbf{3) $\implies$ 4):} Пусть $X \cap Y = X$. Включение 
        $Y \subseteq X \cup Y$ выполняется всегда. Обратно, пусть 
        $x \in X \cup Y$, тогда $x \in X$ или $x \in Y$. Если $x \in X$, 
        то в силу равенства $X = X \cap Y$ имеем $x \in X \cap Y$, откуда 
        $x \in Y$. В обоих случаях $x \in Y$, значит, $X \cup Y \subseteq Y$, 
        и поэтому $X \cup Y = Y$.

        \textbf{4) $\implies$ 1):} Пусть $X \cup Y = Y$. Возьмём любой 
        элемент $x \in X$. По определению объединения $x \in X \cup Y$, но 
        так как $X \cup Y = Y$, получаем $x \in Y$. Следовательно, 
        $X \subseteq Y$.
    \end{quote}
\end{proof}

\subsection{Законы де Моргана для множеств}

\begin{problem}
    Доказать закон де Моргана для дополнения объединения двух множеств:
    \begin{equation} \label{eq:demorgan_cup}
        -(A \cup B) = (-A) \cap (-B) \tag{6}
    \end{equation}
    \emph{Здесь знак унарного минуса обозначает дополнение подмножества до 
    фиксированного универсального множества $U$, то есть $-A = U \setminus A$.}
\end{problem}

\begin{proof}[Решение]~\\
    \begin{quote}
        Докажем равенство множеств \eqref{eq:demorgan_cup} методом 
        взаимного включения.

        \textbf{1. Включение $-(A \cup B) \subseteq (-A) \cap (-B)$.}\\
        Пусть $x \in -(A \cup B)$. По определению дополнения это означает, 
        что $x \in U$ и $x \notin (A \cup B)$. Так как элемент не 
        принадлежит объединению $A \cup B$, он не принадлежит ни одному из 
        множеств $A$ и $B$:
        \[
            x \notin A \quad \text{и} \quad x \notin B
        \]
        По определению дополнения отсюда следует, что $x \in -A$ и 
        $x \in -B$, то есть $x \in (-A) \cap (-B)$.

        \textbf{2. Включение $(-A) \cap (-B) \subseteq -(A \cup B)$.}\\
        Пусть $x \in (-A) \cap (-B)$. По определению пересечения имеем 
        $x \in -A$ и $x \in -B$, то есть $x \notin A$ и $x \notin B$. 
        Следовательно, элемент $x$ не может принадлежать объединению 
        $A \cup B$, откуда $x \notin (A \cup B)$ и $x \in -(A \cup B)$.

        Из двух встречных включений следует равенство 
        $-(A \cup B) = (-A) \cap (-B)$.
    \end{quote}
\end{proof}

\subsection{Метод характеристических функций}

\begin{definition}
    Пусть $U$ --- универсальное множество и $A \subseteq U$. 
    \textdef{Характеристической (индикаторной) функцией} подмножества $A$ 
    называется отображение $\chi_A: U \to \{0, 1\}$, задаваемое правилом:
    \begin{equation} \label{eq:char_func_def}
        \chi_A(x) = 
        \begin{cases}
            1, & \text{если } x \in A, \\
            0, & \text{если } x \notin A.
        \end{cases} \tag{7}
    \end{equation}
\end{definition}

\begin{properties}[Алгебра характеристических функций] ~
    Основные операции над подмножествами $A, B \subseteq U$ однозначно 
    выражаются через арифметические операции над их характеристическими 
    функциями для любого $x \in U$:
    \begin{enumerate}
        \item \textbf{Пересечение:} 
        $\chi_{A \cap B}(x) = \min\bigl(\chi_A(x), \chi_B(x)\bigr) = 
        \chi_A(x) \cdot \chi_B(x)$;
        \item \textbf{Объединение:} 
        $\chi_{A \cup B}(x) = \max\bigl(\chi_A(x), \chi_B(x)\bigr) = 
        \chi_A(x) + \chi_B(x) - \chi_A(x) \chi_B(x)$;
        \item \textbf{Дополнение:} 
        $\chi_{-A}(x) = 1 - \chi_A(x)$;
        \item \textbf{Разность:} 
        $\chi_{A \setminus B}(x) = \chi_A(x) \cdot \bigl(1 - \chi_B(x)\bigr)$;
        \item \textbf{Симметрическая разность:} 
        $\chi_{A \symdiff B}(x) = \bigl(\chi_A(x) + \chi_B(x)\bigr) \bmod 2 = 
        \chi_A(x) + \chi_B(x) - 2\chi_A(x)\chi_B(x)$.
    \end{enumerate}
\end{properties}

\begin{problem}
    Доказать критерий включения множеств на языке характеристических 
    функций:
    \begin{equation} \label{eq:char_sub}
        A \subseteq B \iff \forall x \in U : \chi_A(x) \le \chi_B(x) \tag{8}
    \end{equation}
\end{problem}

\begin{proof}[Решение]~\\
    \begin{quote}
        Докажем равносильность \eqref{eq:char_sub} в обе стороны.

        $\implies$ \textbf{Необходимость.} Пусть $A \subseteq B$ и 
        $x \in U$ --- произвольный элемент. Если $x \notin A$, то 
        $\chi_A(x) = 0 \le \chi_B(x)$, так как $\chi_B(x) \in \{0, 1\}$. 
        Если же $x \in A$, то $\chi_A(x) = 1$, но в силу включения 
        $A \subseteq B$ выполняется $x \in B$, откуда $\chi_B(x) = 1$, и 
        неравенство $1 \le 1$ также верно.

        $\impliedby$ \textbf{Достаточность.} Пусть для всех $x \in U$ 
        выполняется неравенство $\chi_A(x) \le \chi_B(x)$. Возьмём любой 
        элемент $x \in A$, для которого $\chi_A(x) = 1$. Из неравенства 
        $1 \le \chi_B(x) \le 1$ следует $\chi_B(x) = 1$, то есть $x \in B$. 
        Следовательно, $A \subseteq B$.
    \end{quote}
\end{proof}

\begin{problem}
    Доказать с помощью характеристических функций монотонность операции 
    объединения: если $A \subseteq B$, то $A \cup C \subseteq B \cup C$ для 
    любого множества $C$.
\end{problem}

\begin{proof}[Решение]~\\
    \begin{quote}
        По условию $A \subseteq B$, что согласно критерию \eqref{eq:char_sub} 
        равносильно поточечному неравенству $\chi_A(x) \le \chi_B(x)$ для 
        всех $x \in U$.

        Требуется доказать неравенство $\chi_{A \cup C}(x) \le 
        \chi_{B \cup C}(x)$. Воспользуемся представлением характеристической 
        функции объединения через максимум:
        \[
            \chi_{A \cup C}(x) = \max\bigl(\chi_A(x), \chi_C(x)\bigr), \qquad
            \chi_{B \cup C}(x) = \max\bigl(\chi_B(x), \chi_C(x)\bigr).
        \]
        Поскольку функция максимума неубывающая по каждому аргументу, из 
        неравенства $\chi_A(x) \le \chi_B(x)$ немедленно следует:
        \[
            \max\bigl(\chi_A(x), \chi_C(x)\bigr) \le 
            \max\bigl(\chi_B(x), \chi_C(x)\bigr) \implies 
            \chi_{A \cup C}(x) \le \chi_{B \cup C}(x).
        \]
        По критерию \eqref{eq:char_sub} это означает, что 
        $A \cup C \subseteq B \cup C$.
    \end{quote}
\end{proof}

\subsection{Мощность множеств и построение биекций}

\begin{definition}
    Два множества $A$ и $B$ называются \textdef{равномощными} (обозначается 
    $\abs{A} = \abs{B}$ или $A \sim B$), если между ними существует 
    биективное отображение $f: A \to B$. Множество называется 
    \textdef{счётным}, если оно равномощно множеству натуральных чисел $\N$.
\end{definition}

\begin{problem}
    Доказать счётность множества всех целых чисел $\Z$ и множества всех 
    многочленов с целыми коэффициентами $\Z[x]$.
\end{problem}

\begin{proof}[Решение]~\\
    \begin{quote}
        \textbf{1. Счётность множества $\Z$.} 
        Упорядочим элементы множества $\Z$ «чередованием» знаков: 
        $0, 1, -1, 2, -2, 3, -3, \dots$. Зададим биекцию $f: \N \to \Z$ 
        явной аналитической формулой (считая $\N = \{1, 2, 3, \dots\}$):
        \begin{equation} \label{eq:bij_z}
            f(n) = 
            \begin{cases}
                \dfrac{n}{2}, & \text{если } n \text{ чётно}, \\[2mm]
                -\dfrac{n - 1}{2}, & \text{если } n \text{ нечётно}.
            \end{cases} \tag{9}
        \end{equation}
        Чётные натуральные числа биективно переходят во все натуральные 
        положительные числа, а нечётные --- в ноль и все целые 
        отрицательные числа. Следовательно, $\Z \sim \N$.

        \textbf{2. Счётность кольца многочленов $\Z[x]$.} 
        Любой многочлен степени не выше $k$ однозначно определяется 
        кортежем длины $k + 1$ своих целых коэффициентов:
        \[
            P(x) = a_0 + a_1 x + \dots + a_k x^k \longleftrightarrow 
            (a_0, a_1, \dots, a_k) \in \Z^{k+1}
        \]
        Так как декартова степень $\Z^{k+1}$ счётного множества $\Z$ при 
        любом фиксированном $k \in \N \cup \{0\}$ счётна, то и множество 
        $M_k$ многочленов степени не выше $k$ счётно. Всё множество 
        многочленов $\Z[x] = \bigcup_{k=0}^\infty M_k$ является 
        объединением счётного семейства счётных множеств, а значит, счётно.
    \end{quote}
\end{proof}

\begin{problem}
    Доказать, что замкнутый отрезок $[0, 1]$ равномощен всей вещественной 
    прямой $\R$.
\end{problem}

\begin{proof}[Решение]~\\
    \begin{quote}
        Построим искомую биекцию в два этапа, установив сначала 
        равномощность $(0, 1) \sim \R$, а затем $[0, 1] \sim (0, 1)$.

        \textbf{1. Биекция между открытым интервалом $(0, 1)$ и прямой $\R$.}\\
        Рассмотрим масштабированную и сдвинутую функцию тангенса 
        $f: (0, 1) \to \R$:
        \begin{equation} \label{eq:bij_tan}
            f(x) = \tg\left(\pi x - \frac{\pi}{2}\right) \tag{10}
        \end{equation}
        На интервале $x \in (0, 1)$ аргумент $\left(\pi x - \frac{\pi}{2}\right)$ 
        пробегает интервал $\left(-\frac{\pi}{2}, \frac{\pi}{2}\right)$, на 
        котором функция тангенса строго возрастает, непрерывна и принимает 
        все вещественные значения от $-\infty$ до $+\infty$. Следовательно, 
        отображение \eqref{eq:bij_tan} биективно.

        \textbf{2. Биекция между отрезком $[0, 1]$ и интервалом $(0, 1)$.}\\
        Применим конструкцию «отеля Гильберта». Выделим внутри интервала 
        $(0, 1)$ счётную последовательность попарно различных точек 
        $a_n = \frac{1}{n}$ при $n = 2, 3, 4, \dots$ 
        ($\frac{1}{2}, \frac{1}{3}, \frac{1}{4}, \dots$). «Освободим» первые 
        два места этой последовательности для граничных точек $0$ и $1$, 
        сдвинув все элементы последовательности на два номера вперёд. 
        Зададим отображение $g: [0, 1] \to (0, 1)$ формулой:
        \begin{equation} \label{eq:bij_hilbert}
            g(x) = 
            \begin{cases}
                \dfrac{1}{2}, & \text{если } x = 0, \\[2mm]
                \dfrac{1}{3}, & \text{если } x = 1, \\[2mm]
                \dfrac{1}{n + 2}, & \text{если } x = \dfrac{1}{n} 
                \text{ при } n \in \N, \ n \ge 2, \\[2mm]
                x, & \text{для всех остальных } x \in (0, 1).
            \end{cases} \tag{11}
        \end{equation}
        Построенное отображение $g$ инъективно и сюръективно на $(0, 1)$. 
        Тогда композиция $h = f \circ g: [0, 1] \to \R$ двух биекций 
        \eqref{eq:bij_tan} и \eqref{eq:bij_hilbert} также является 
        биекцией, что доказывает равномощность $[0, 1] \sim \R$.
    \end{quote}
\end{proof}